%\chapter{Reference guide}

\section{Additional packages available with Astra--7.0}
%
\subsection{TORBEAM for ECRH/ECCD calculations}
Recently, the TORBEAM code has been added to ASTRA through an
interface subroutine \textit{torba.f} that is found in the sbr/ directory. 

The executable files are to be extracted from the library {\it libtorbeam.a}, 
and the resulting {\it xxx.o} are to be added to the {\it user.a} archive of 
the Astra user. 

To call TORBEAM just call it as a subroutine in the model file as:
\\

TORBA(VAR1,VAR2)::::;
\\

The outputs {\bf VAR1,VAR2} are two arrays. {\bf VAR1} is the ECRH power density
deposition profile in [MW/m$^3$], while {\bf VAR2} is the ECCD profile in 
[MA/m$^2$]. Note that the power profile is sumed up over all gyrotrons, as well 
as the driven current profile.

At present TORBA can run up to 8 gyrotrons. Their specifications are to be 
prescribed in a file that has the same name as the experimental file, and has to
be placed in the {\it exp/ecr/} directory. The structure of the ECRH file is the
same as the NBI file explained in section 4, table 4.30. In the following table 
6.1.1 the different parameters are explained for the ECRH file:\\
\noindent
{\bf Table 6.1.1. Parameters describing each gyrotron}\\[2ex]
\begin{tabular}{| c | c | l |}
\hline
\bf No.&\bf\hspace*{-1pt}Units\hspace*{-1pt}&\bf Description\\
\hline
1  &[MW]& Gyrotron nominal power					\\
2  &--& Unused parameter \\
3  &[GHz]& Frequency\\
4  &[cm]& Major radius from where beam is launched	\\
5  &[cm]& Vertical position from where beam is launched	\\
6  &[deg]& Angle of poloidal injection (from LFS, positive for upward launching) 	\\
7  &[deg]& Angle of toroidal injection (from LFS, positive for clockwise launching)	\\
8--20 &--& Unused parameters	\\
\hline
\end{tabular}\\[15pt]
The first parameter gives the power injected by each gyrotron.
This quantity is usually time dependent, therefore, it is convenient to
set it as a variable rather than as a number. 
For example, this variable can be a $\tt ZRD$
and its time evolution can be prescribed in the data file.
%
\subsection{Quasi--linear transport model TGLF}
The TGLF code [G. Staebler {\it et al.}, Phys. Plasmas {\bf 14}, 055909 (2007)] 
computes turbulence--driven fluxes for heat and particle transport via solution
of fluid moments with kinetic closure fitted to match kinetic responses computed
with linear gyrokinetic codes. The trapped and passing particle descriptions 
have been majorly upgraded with respect to its predecessor GLF23.

In principle different interfaces between Astra and TGLF are available, for both
serial and parallel computations. They should be requested explicitly.
%
\subsection{Quasi--linear transport model QuaLiKiz}
%The TGLF code [G. Staebler {\it et al.}, Phys. Plasmas {\bf 14}, 055909 (2007)] 
%computes turbulence--driven fluxes for heat and particle transport via solution
%of fluid moments with kinetic closure fitted to match kinetic responses computed
%with linear gyrokinetic codes. The trapped and passing particle descriptions 
%have been majorly upgraded with respect to its predecessor GLF23.

%In principle different interfaces between Astra and TGLF are available, for both
%serial and parallel computations. They should be requested explicitly.
%
\subsection{TORIC for ICRF calculations}
An interface to couple Astra to the TORIC code for ICRF heating computations has
been developed. The routine is called {\it a2toric.f} and is part of the {\it sbr/} package. The TORIC library however has to be requested explicitly.
%
\subsection{STRAHL for impurity transport calculations}
Astra is also now coupled to the impurity transport code STRAHL, such that 
time--dependent computation can be performed with STRAHL providing impurity 
profiles, $Z_{\rm eff}$ profiles, and radiated power profiles.

The interface is called {\it A2STRAHL.f} and is part of the {\it sbr/} package. 
The STRAHL code itself has to be requested explicitly. 
%
\subsection{NEO for neoclassical calculations}
%Astra is also now coupled to the impurity transport code STRAHL, such that 
%time--dependent computation can be performed with STRAHL providing impurity 
%profiles, $Z_{\rm eff}$ profiles, and radiated power profiles.

%The interface is called {\it A2STRAHL.f} and is part of the {\it sbr/} package. 
%The STRAHL code itself has to be requested explicitly. 
%
\subsection{ASTRA restart files}
Sometimes ASTRA computations can be lengthy, and it is important to be able to 
save relevant data such that one can restart the simulation exactly as it has 
been saved.

This new routine is called {\it restart\_astra.f} and is in the {\it sbr/} 
package. 
It can be called in the model file as:

RESTARTA::::;
\\
as with all the other routines, it can be assigned a starting time, an ending 
time, and time intervals for execution, or keys to be pushed to activate it.
Each time the routine is called, it will create a new subdirectory in the {\it 
dat/} directory, the subdirectory being called {\it 
(model)\_(expfile)\_(index)}, where (model) is the model file name, (expfile) 
is the experimental file name, and (index) is just an index ordering the time 
slices.

Also, a new control parameter has been added, which is called {\bf IRESTA}, that can have the following values:
\\
{\bf Table 6.5.1} Choice of equilibrium solution with {\bf IRESTA}\\[2ex]
\begin{tabular}{|l|l|l|}					\hline
\bf Name	&\bf Value & Description\\ \hline
\tt IRESTA 	& 0 & Do not use any restart file\\
\tt  	& $i$ (integer $>$ 0) &  Use restart file with index $i$\\ 
\hline
\end{tabular}
\bigskip\\
By default, {\bf IRESTA}=0.

If {\tt IRESTA}$>$0, then the restart file with index {\tt IRESTA} is used. What
the restart option does essentially is reading first the experimental file and 
model log file as they are, then reading the restart file and substituting all 
constants with the values as they are stored in the restart file. The constants 
include all variables and control/grid parameters. The simulation is then 
started from the {\tt TIME} as stored in the restart file.
As regards profiles, contrary to the usual Astra initialization, when using the restart 
file, the following profiles are substituted as they are stored in the restart 
file:

{\tt TE, NE, TI, F(1--12), FP, UPAR, MU, CU}\\
Instead of using the values as given in the exp file. This allows to restart a simulation exactly from the state in which it was left before.
%
\clearpage
